% ===== 原文 pp.87--98 =====
% tr07a.tex —— §15 高阶代数 K 理论（pp.87--88）、§16 第一讲（pp.89--98）
% 译者：tr-07。展示式一律 \[ … \]，定理类号数照原文、每节从 1 起。

% [原文 p.87]
\section{高阶代数 K 理论}

代数 K 理论对一切整数 $n$ 定义并应用一列记作 $K_n$ 的函子，它们从环出发、取值于阿贝尔群。环既可以是交换的，也可以是非交换的；通常是含幺环，但后一假设可以放松，代价是复杂度增加。

% [Tutor Note: p.87 末段 "with his trully ingeneous plus construction" 有两处拼写错误：trully 应为 truly、ingeneous 应为 ingenious（500 dpi 裁图复核 p.87 末行）；同段 "this later assumption" 亦应为 "this latter assumption"（指前一句的 unital）。译文照原样处理，不给原文改字。]
由于历史原因，人们设想低阶代数 K 群 $K_0$ 与 $K_1$ 的方式，与设想高阶代数 K 群 $K_n$（$n \ge 2$）的方式有所不同。的确，低阶群比高阶群更容易接近，并且（直到今天）有更多的应用。高阶 K 群的理论明显更为深刻，而且计算起来肯定困难得多（即使所取的环是整数环）。

群 $K_0$ 借助投射模推广了环的理想类群的构造（后者起源于整数唯一分解性质的失效）。它在 1960 年代与 1970 年代的发展，与试图解决 Serre 关于投射模的一个猜想相联系——这个猜想如今就是 Quillen--Suslin 定理；在那个时代还发现了它与经典代数问题的许多其他联系。类似地，$K_1$ 借助初等矩阵理论，是对环中单位群的一种修正。直观地说，可以认为 $K_1$ 编码了以环中元素为元素的 $\infty \times \infty$ 矩阵的同伦类的信息。

第一个 K 群 $K_1(R)$ 在拓扑中很重要，尤其当 $R$ 是群环时，因为它的商——Whitehead 群——包含用来研究简单同伦论与手术理论问题的 Whitehead 挠；群 $K_0(R)$ 还包含其他不变量，例如有限性不变量。自 1980 年代以来，代数 K 理论在代数几何中的应用越来越多。例如，动机上同调与代数 K 理论密切相关。

第 0 个 K 群由 Alexander Grothendieck 在约 1957 年研究代数几何时定义；此后不久，第 1 个 K 群由著名拓扑学家 John Whitehead 定义；而第 2 个 K 群由 Robert Steinberg 定义。进展一直停滞，直到 1970 年代 D.G. Quillen 出现；凭借他那极其精巧的\emph{+ 构造}，他给出了第 $n$ 个 K 群的一般定义，其中包含 Grothendieck、Whitehead 与 Steinberg 对 $n = 0, 1, 2$ 的全部先前定义。

% [原文 p.88]
若干年后，Quillen 发现了他著名的 $Q$-构造；它更为一般，函子性质有所改进，并且可以应用于任何正合范畴（而不只是环）。Quillen 的结果使 Grothendieck 大为振奋，后者寄出了一封 600 页的手写信，论及代数几何与 K 理论的未来。在 1980 年代，Quillen 高阶代数 K 理论的一个拓扑版本由 Waldhausen 发展出来，称为 Waldhausen K 理论。

% [Tutor Note: p.88 末句 "the computations are only known ”mudulo torsion” (Borel)" 中 mudulo 应为 modulo（引号方向在原文里也是反向的）；所指即 Borel 对整数环 $K_n(\Z)$ 的计算：秩已知，挠部分未知。]
不幸的是，正如现代数学中通常的情形，计算很困难，因此即使在今天也只有少数已知的例子：完全的计算只对有限域存在（Quillen）；而即使对整数环，计算也只是在「模去挠」的意义下已知（Borel）。

% [原文 p.89]
\section{第一讲（代数预备与 Grothendieck 群）}

% [Tutor Note: p.89 §16 标题原文作 "Lecture 1 (Algebraic Preliminaries and the the Grothendieck Group)"，重复了一个 "the"（页图上第二行开头即 "the Grothendieck Group"）；同页 "abreviated" 应为 abbreviated，Definition 1 中 "projctive" 应为 projective（均经 500 dpi 复核）。标题按 BRIEF 规定的中文写法给出，正文照原样翻译。]
本章的关键概念将是 $A$-模的概念，其中 $A$ 是一个环（交换或非交换，通常但不必定含幺）。当 $A$ 实际上是一个域时，$A$-模就只是一个 $A$-向量空间。受向量空间研究中通常记号的启发，我们将考虑左 $A$-模。

线性代数中的经典结果（有限维向量空间）说：每个向量空间都有基，且某个固定向量空间的所有基中向量的个数是常数；这使人们能够定义向量空间的维数这一概念。对于模，类似（但不完全相同）的事情也成立：$A$-模 $P$ 的一组基是 $P$ 的一个子集 $\{e_i\}_{i \in I}$（不必具有有限基数；若如此，则称 $P$ 为有限生成的，简记为 f.g.，见下文），使得 $P$ 的每个元素都能唯一地写成有限和 $\sum_i a_i e_i$，其中 $a_i \in A$。若 $P$ 有一组固定的基，则称它为有基自由模，并把 $P$ 的维数（或秩）定义为其给定基的基数。若存在一组基使模成为有基自由模，则称该模为自由的。自由模的典型例子是 $A^n$，它由环 $A$ 的元素的 $n$ 元组组成。然而有些环满足 $A^n \simeq A^{n+m}$，其中 $m \ne 0$。为避免这种病态，我们将假定我们所有的环都具有所谓的基不变性质，其确切含义是：除非 $n = m$，$A^n$ 与 $A^m$ 不同构。在这种情形下，自由 $A$-模 $P$ 的维数是不变的，与 $P$ 中基的选择无关。（一个不具有基不变性质的模的例子是无限矩阵环 $\End_F(F^\infty)$，即某个域 $F$ 上无限维向量空间 $F^\infty$ 的自同态的环。）在简短的代数预备之后，我们进入本节的基本材料。

% [Tutor Note: p.89 "（一个不具有基不变性质的模的例子是无限矩阵环 $\End_F(F^\infty)$）" 原文说它是 "a module"，其实它是环：把 $F^\infty \simeq F^\infty \oplus F^\infty$ 搬过去即得 $\End_F(F^\infty) \simeq \End_F(F^\infty)^2$，故这个环的自由模破坏基不变性质。]
\begin{defn}{1} 设 $A$ 是含幺环。一个 $A$-模（$A$-module）$P$ 称为\emph{投射的}，若存在一个模 $Q$，使得直和项
\[ F = P \oplus Q \]
是自由模 $F$。\end{defn}

% [原文 p.90]
由上述定义可知，$P$ 是 $F$ 上某个投影 $p$ 的像；$F$ 中那个在 $P$ 上为恒等、在 $Q$ 上为 $0$ 的模自同态，即
\[ Q = \Ker(p), \]
是幂等的，并把 $F$ 投影到 $P$ 上。换言之，$P$ 是投射的，当且仅当 $P$ 是自由模的直和项。

用范畴的语言，可以借助提升性质来定义投射模：一个 $A$-模 $P$ 是投射的，当且仅当对每个满射模同态 $f : N \onto M$ 和每个模同态 $g : P \to M$，都存在同态 $h : P \to N$ 使得 $fh = g$（我们不要求提升同态 $h$ 唯一，这不是泛性质）。下图描述这一情形：

% [Tutor Note: p.90 下图左竖线是恒等（原文用 $\|$），右竖线是 $f : N \to M$；作者特意声明 h 不唯一，故这只是一条「提升性质」而非泛性质——泛性质要求相应的映射唯一。]
\[
\begin{tikzcd}
P \arrow[r,"h"] \arrow[d,equal] & N \arrow[d,"f"] \\
P \arrow[r,"g"] & M
\end{tikzcd}
\]

有限生成模的概念推广了有限维向量空间的概念（当环 $A$ 实际上是域时）：

\begin{defn}{2} 一个（左）$A$-模 $P$ 称为\emph{有限生成的}，当且仅当存在有限多个元素 $x_1, x_2, \dots, x_n \in P$，使得 $\forall x \in P$，都存在元素 $a_1, a_2, \dots, a_n \in A$，使
\[ x = a_1x_1 + a_2x_2 + \dots + a_nx_n. \]
集合 $\{x_1, x_2, \dots, x_n\}$ 称为 $P$ 的一组生成元。\end{defn}

\begin{thm}{1} 以下条件等价：
\begin{enumerate}[label=\arabic*.]
\item $P$ 是有限生成投射（简记 f.g.p.）$A$-模。
\item 存在 $n \in \N^*$ 以及另一个 $A$-模 $Q$，使得
\[ P \oplus Q = A^n. \]
% [原文 p.91]
\item 存在某个 $n \in \N^*$ 以及某个元素 $e \in M_nA$，使得
\[ e = e^2 \]
且
\[ P = A^n e \]
（其中 $M_nA$ 表示元素取自 $A$ 的 $n \times n$ 矩阵的集合）。
\end{enumerate}
\end{thm}

% [Tutor Note: p.90 定理 1 的 (2)$\Rightarrow$(3)：取 $e$ 为沿 $Q$ 的投影矩阵，则 $P = A^ne$；反向由 $e^2 = e$ 得 $A^n = \operatorname{Im}e \oplus \operatorname{Ker}e$，两个直和项都是投射模。原文把证明推给引言所列参考书，故此处只列等价条件。]
本定理的证明请参看引言中列出的代数参考书（背景文献）。

我们用 $P_A$ 表示有限生成投射 $A$-模的范畴，用 $\operatorname{Iso}(P_A)$ 表示有限生成投射 $A$-模的同构类的集合。第二个集合实际上是一个阿贝尔幺半群，其加法由
\[ [P] + [Q] = [P + Q]. \]
给出。

% [Tutor Note: p.91 定义 3 在 $\operatorname{Iso}(P_A)$ 之后多印了一个右括号（页图为 "Iso(P_A) )"），译文照原样保留；同段 "the Grothendieck group $Gr$ we used in chapter 1" 指第 1 章由 $\Vect(X)$ 构造 Grothendieck 群的同一手法。]
\begin{defn}{3} 我们定义环 $A$ 的第 0 个 K 群，记作 $K_0A$，即 $\operatorname{Iso}(P_A)$ 的阿贝尔群完备化（或我们在第 1 章用过的 Grothendieck 群 $\operatorname{Gr}$）），即
\[ K_0A := \operatorname{Iso}(P_A)\ \text{的阿贝尔群完备化}. \]
\end{defn}

第 0 个 K 群常被称为环 $A$ 的\emph{Grothendieck 群}，以纪念伟大的法国数学家（德裔）Alexander Grothendieck，他在 1957 年定义了它。

一般地，若 $I$ 是任意阿贝尔幺半群，则存在一个阿贝尔群 $I^*$ 连同同态 $\phi : I \to I^*$，它是从 $I$ 到某个阿贝尔群的泛同态。$(I^*, \phi)$ 只在规范同构的意义下确定。

构造 $I^*$ 有三种方式：
\begin{enumerate}[label=\arabic*.]
\item $I^*$ 是以 $[a]$（$a \in I$）为生成元、以
\[ [a + b] = [a] + [b]. \]
为关系的自由阿贝尔群。

% [原文 p.92]
\item $I^* = I \times I/\sim$，其中等价关系 $\sim$ 定义如下：$(a, b) \sim (a', b') \iff \exists\gamma : a + b' + \gamma = a' + b + \gamma$（显然这确实是一个等价关系）。其次我们定义 $(a, b) + (a', b') = (a+a', b+b')$，而 $0 = (0, 0)$、$-(a, b) = (b, a)$。于是 $(a, b) = [a] - [b]$。（这正是我们在第 1 章 1.1 节中用来定义 $K(X)$ 的构造，即 $\Vect(X)$ 的 Grothendieck 群。）

\item 假定存在 $a_0 \in I$，使得 $\forall a' \in I$，$\exists n \in \N$ 与 $b \in I$ 满足 $a + b = nd_0$。则
\[ I^* = I \times \N/\sim \]
其中 $(a, n) \sim (a', n') \iff \exists n \in \N$ 使得 $a + n'a_0 = a' + na_0$。
\end{enumerate}

% [Tutor Note: p.92 第 3 种构造里 "a + b = nd_0" 的 d_0 应为 a_0（同一段下一行两次写作 $a_0$，页图 500 dpi 复核）；此外 "$\forall a' \in I$" 之前原文漏掉了对 $a$ 的量化，正确的陈述应是：存在 $a_0 \in I$，使 $\forall a \in I$ 都有 $b, n$ 满足 $a + b = na_0$。]
以上所有阿贝尔群都具有所要求的泛性质。

于是，总结起来，含幺环 $A$ 的第 0 个 K 群 $K_0A$ 是以 $[p]$（$\forall p \in P_A$）为生成元、以 $[p \oplus q] = [p] + [q]$ 为关系的自由阿贝尔群，它等于形式差 $[p] - [q]$ 构成的群，也等于形式差 $[p] - [A^n]$ 构成的群。

\para{例：}
\begin{enumerate}[label=\arabic*.]
\item 设 $\F$ 是除环（skew-field），即满足 $0 \ne 1$、其中除法可行的环；换言之，每个元素都有乘法逆元（等价地说，是一个乘法不必交换的域）。那么有限生成投射 $\F$-模的范畴 $P_{\F}$ 由 $\F$ 上的有限维向量空间组成。于是阿贝尔幺半群 $\operatorname{Iso}(P_{\F})$ 等于 $\N$，从而
\[ K_0\F = \Z. \]

\item 若 $A = \Z$，则对一切 $n \in \N$，范畴 $P^n_{\Z}$ 就是有限生成自由阿贝尔群 $\Z^n$，而
\[ \operatorname{Iso}(P_{\Z}) = \N, \]
从而
\[ K_0\Z = \Z. \]

% [Tutor Note: p.92 例 2 的 "the category $P^n_{\Z}$ is the finitely generated free abelian group $\Z^n$" 一句在原文即不通（范畴不可能是群）；按上下文应读作：$P_{\Z}$ 中的有限生成投射对象就是 $\Z^n$（$n \in \N$）。译文照原文，未经改动。]
% [原文 p.93]
主理想整环也一样。

\item Dedekind 整环，例如取 $\F$ 为某个数域 $[F, Q] < \infty$，并取 $A$ 为 $\Z$ 在 $\F$ 中的整商。以 $\operatorname{Pic}(A)$ 表示 $A$ 的分式理想的理想类群除以其主理想所得的群。于是有下述结论：

% [Tutor Note: p.93 这里的 "Theorem 2:" 原文用斜体加冒号排版，与该节其余定理的粗体体例不一致；同句 "the integral quotient of $\Z$ in $\F$" 应为 "integral closure"（整闭包），"$[F, Q]$" 的通常写法是 $[F:Q]$（$\F$ 在 $\Q$ 上的次数），末句 "excercise" 应为 exercise。以上均照译。]
\begin{thm}{2} 若 $P \in P_A$，则 $P \simeq a_1 \oplus a_2 \oplus \dots \oplus a_n$，其中 $a_i$ 是分式理想。因此 $K_0A = \Z \oplus \operatorname{Pic}(A)$。\end{thm}

本定理的证明留给读者作为练习（例如见 [39]）。

\item \textbf{Serre--Swan 定理.} 设 $X$ 是紧 Hausdorff 拓扑空间，$A = C(X)$ 表示 $X$ 上连续的复值函数的集合（它是一个含幺环）。那么范畴 $P_A$ 等价于 $X$ 上复向量丛的范畴。

% [Tutor Note: p.93 注意 $A = C(X)$ 取的是复值连续函数，故对偶的是复向量丛；若取实值函数则对应实向量丛。本条（列表第 4 项）不算编号定理，原文只是把它排成列表第 4 项且标题加粗。]
\end{enumerate}

\emph{［旁注：本定理是法国数学家 Alain Connes 发展非交换几何的动因之一；非交换几何在某种意义上试图统一微分几何与代数几何。］}

\para{Serre--Swan 定理的证明：}
这是一个重要的定理，因此我们将给出它的完整证明。我们必须证明 $\operatorname{Iso}(P_A)$（即有限生成投射 $A$-模的同构类的集合）与 $\Vect(X)$（即 $X$ 上复向量丛的同构类的集合）之间存在双射，其中 $A = C(X)$。

回忆第 1 章第 1 节中向量丛与向量丛的局部截面的定义。给定任意向量丛 $E \to X$，我们用 $\Gamma(E, X)$（或简记 $\Gamma(E)$）表示向量丛 $E \to X$ 的所有局部截面构成的集合。显然，若定义：
\begin{itemize}
\item $(s_1 + s_2)(x) = s_1(x) + s_2(x)\forall s_1, s_2 \in \Gamma(E)$
\item $(as)(x) = a(x)s(x)\forall a \in C(X)$ 与 $s \in \Gamma(E)$
\end{itemize}
则 $\Gamma(E)$ 成为一个 $C(X)$-模。

显然，对平凡向量丛 $\C^n_X$，集合 $\Gamma(\C^n_X)$ 是 $n$ 个生成元上的自由 $C(X)$-模。此外，$\Gamma(-)$ 是从 $X$ 上复向量丛的范畴到范畴 $P_A$ 的可加函子。因此，为证明本定理，我们必须证明三个断言：

% [原文 p.94]
\begin{claim}{1} 函子 $\Gamma(-)$「保持」同构，即若 $E \simeq E'$ 作为 $X$ 上的向量丛同构，则 $\Gamma(E) \simeq \Gamma(E')$ 作为 $C(X)$-模同构。\end{claim}

\begin{claim}{2} 对 $X$ 上任意向量丛 $E$，$\Gamma(E)$ 是有限生成投射 $C(X)$-模。\end{claim}

\begin{claim}{3} 所有有限生成投射 $C(X)$-模都可「实现」为某个 $X$ 上向量丛 $E$ 的 $\Gamma(E)$。\end{claim}

断言 1 的证明需要若干命题（我们在更一般的情形下证明它，即 $X$ 是正规拓扑空间；更确切地说，我们将证明：若 $X$ 正规，则 $\Gamma(-)$ 给出同构 $\Hom(E, E') \simeq \Hom_{C(X)}(\Gamma(E), \Gamma(E')))$）：

% [Tutor Note: p.94 中段那个 Hom 同构句末多印了一个右括号（页图为 "$\Gamma(E')))$"），译文照原样保留；两处多余括号（p.91 定义 3 与这里）是同一类排版笔误。]
\begin{lem}{1} 若 $X$ 正规，设 $U$ 是 $x \in X$ 的一个邻域，$s$ 是向量丛 $E$ 在 $U$ 上的一个截面。则存在 $E$ 在 $X$ 上的截面 $s'$，使得 $s$ 与 $s'$ 在 $x$ 的某个邻域上一致。\end{lem}

% [Tutor Note: p.94 引理 1 的证明真正用到的是 $X$ 的正规性：$\bar W$ 是 $x$ 的闭邻域、$X - V$ 是闭集，二者不交，由 Urysohn 引理可找到连续函数 $a$ 在 $\bar W$ 上等于 1、在 $X - V$ 上等于 0。取 $W \subset \bar W \subset V \subset \bar V \subset U$ 的嵌套是使 $s'$ 在 $U$ 上良定义且延拓为 0 的关键。]
\begin{proof} 我们将构造 $s'$：取 $x$ 的邻域 $V, W$，使得 $\bar V \subset U$ 且 $\bar W \subset V$。设 $a$ 是 $X$ 上的实值函数，满足 $a|_{\bar W} = 1$ 且 $a|_{(X-V)} = 0$。我们令 $s'(y) = a(y)s(y)$（若 $y \in U$）以及 $s'(y) = 0$（若 $y \notin U$）。\end{proof}

\begin{cor}{1} 若 $X$ 正规，则对任意 $x \in X$，存在元素 $s_1, s_2, \dots, s_n \in \Gamma(E)$，它们在 $x$ 处构成一组局部基。\end{cor}

\begin{proof} （显然。）\end{proof}

\begin{cor}{2} 若 $X$ 正规，$f, g : E \to E'$ 是两个向量丛映射，且 $\Gamma(f) = \Gamma(g) : \Gamma(E) \to \Gamma(E')$，则 $f = g$。\end{cor}

\begin{proof} 给定 $e \in E$，$\pi(e) = x$（其中 $\pi : E \to X$ 是向量丛投影），存在 $x$ 的某个邻域 $U$ 上的截面 $s$，满足 $s(x) = e$。由引理 1，存在截面 $s' \in \Gamma(E)$ 使 $s'(x) = e$。于是 $f(e) = fs'(x) = (\Gamma(f)s')(x) = (\Gamma(g)s')(x) = g(e)$。\end{proof}

\begin{lem}{2} 若 $X$ 正规，且 $s \in \Gamma(E)$ 满足 $s(x) = 0$，则存在
% [原文 p.95]
元素 $s_1, \dots, s_k \in \Gamma(E)$ 与 $a_1, \dots, a_k \in C(X)$，使得对 $i = 1, \dots, k$ 有 $a_i(x) = 0$，且 $s = \sum_i a_i s_i$。\end{lem}

% [Tutor Note: p.95 引理 2 证明中 "vanishes in an neighborhood" 语法有误（应为 a neighborhood）。关键等式 $s = as' + \sum_i a_i s_i$ 要分两段验证：在 $U$ 内 $s' = 0$，两边都等于 $\sum_i a_i s_i$；在 $X - U$ 上 $a = 1$，两边都等于 $s' + \sum_i a_i s_i = s$。这是作者省略的论证。]
\begin{proof} 由推论 1，设 $s_1, \dots, s_n \in \Gamma(E)$ 是 $x$ 处的一组局部基，并在 $x$ 附近有 $s(y) = \sum_i b_i(y)s_i(y)$，其中 $b_i(y) \in \C$。取 $a_i \in C(X)$ 使 $a_i$ 与 $b_i$ 在 $x$ 的某个邻域上一致（把引理 1 应用于 $X \times \C$，这样的 $a_i$ 存在）。于是 $s' = s - \sum_i a_i s_i$ 在 $x$ 的某个邻域 $U$ 上为零。取 $x$ 的邻域 $V$ 使 $\bar V \subset U$。取 $a \in C(X)$ 在 $x$ 处为零、在 $X - V$ 上为 1。则 $s = as' + \sum_i a_i s_i$。但 $a(x) = 0$ 且 $a_i(x) = b_i(x) = 0$。\end{proof}

\begin{cor}{3} 设 $I_x$ 是 $C(X)$ 中由所有满足 $a(x) = 0$ 的 $a \in C(X)$ 组成的双边理想。则 $\Gamma(E)/[I_x\Gamma(E)] \simeq \pi^{-1}(x)$（其中 $\pi : E \to X$ 是向量丛投影），该同构由 $s \mapsto s(x)$ 给出。\end{cor}

\begin{proof} 这由引理 2 与推论 2（的证明）推出。\end{proof}

\begin{prop}{1} 若 $X$ 正规，则给定任意 $C(X)$-映射 $F : \Gamma(E) \to \Gamma(E')$，存在唯一的（复）向量丛映射 $f : E \to E'$，使得 $F = \Gamma(f)$。\end{prop}

\begin{proof} 唯一性由推论 2 推出。现在 $F$ 诱导出映射 $f_x : \Gamma(E)/[I_x\Gamma(E)] \to \Gamma(E')/[I_x\Gamma(E')]$。所有这些映射合起来给出映射 $f : E \to E'$，它在纤维上是线性的。若 $s \in \Gamma(E)$，则按构造有 $(fs)(x) = f_x s(x) = (F(s))(x)$，从而 $F = \Gamma(f)$。最后一步是验证连续性：设 $s_1, \dots, s_m \in \Gamma(E)$ 是 $x$ 处的一组局部基。若 $e \in E$ 且 $\pi(e)$ 在 $x$ 附近，则有 $e = \sum_i a_i(e)s_i(\pi(e))$，其中 $a_i$ 是连续复值函数。现在 $f(e) = \sum_i a_i(e)fs_i(\pi(e))$。由于 $fs_i = F(s_i)$，这说明 $fs_i$ 是 $E'$ 的连续截面；而和式中所有项都连续依赖于 $e$，故 $f$ 连续。\end{proof}

\begin{cor}{4} 设 $X$ 正规，$E$ 与 $E'$ 是 $X$ 上的两个复向量丛。则 $E \simeq E' \iff \Gamma(E) \simeq \Gamma(E')$ 作为 $C(X)$-模。\end{cor}

% [Tutor Note: p.95 推论 4 原文作 "let $E$ and $E'$ be to complex vector bundles"，其中 to 应为 two；同页命题 1 陈述中的 "given any $C(X)$-map" 指 $C(X)$-模同态。]
\begin{proof} 由命题 1 显然。\end{proof}

为证明断言 2，我们需要两个引理：

% [原文 p.96]
\begin{lem}{3} 若 $X$ 紧且 Hausdorff，设 $E$ 是 $X$ 上任意复向量丛。则存在平凡复向量丛 $\C^n_X$（对某个 $n \in \N$）连同满射 $f : \C^n_X \to E$。\end{lem}

% [Tutor Note: p.96 引理 3 证明引用的 "Propositon 1.1.1"（拼写应为 Proposition）即第 1 章关于紧 Hausdorff 底空间上每个复向量丛都有补丛的命题；其证明用单位分解与 Hermite 度量。作者说明只在实情形证过、复情形可照搬，这里引用时未复述。]
\begin{proof} 回忆命题 1.1.1 断言：若底空间 $X$ 紧且 Hausdorff，则对 $X$ 上每个复向量丛 $E$，都存在 $X$ 上另一个复向量丛 $E'$，使它们的直和是平凡丛，即
\[ E \oplus E' \simeq \C^n_X \]
对某个 $n \in \N$。（我们对实情形证明过该定理，但可以直截了当地搬到复情形。）于是引理 3 是该命题的直接推论，其中满射 $f$ 就是自然投影
\[ f : \C^n_X \simeq E \oplus E' \to E. \]
\end{proof}

\begin{lem}{4} 对 $X$ 上任意复向量丛 $E \to X$（其中 $X$ 紧且 Hausdorff），$\Gamma(E)$ 是有限生成投射 $C(X)$-模。\end{lem}

\begin{proof} 显然，对平凡向量丛 $\C^n_X$，集合 $\Gamma(\C^n_X)$ 是 $n$ 个生成元上的自由 $C(X)$-模。由于 $E \oplus E' \simeq \C^n_X$，由推论 4 便得 $\Gamma(E) \oplus \Gamma(E') \simeq \Gamma(\C^n_X)$ 作为 $C(X)$-模，因此 $\Gamma(E)$ 是有限生成投射 $C(X)$-模。\end{proof}

最后，为证明断言 3，我们需要下述引理与命题：

\begin{lem}{5} 设 $t_1, \dots, t_k$ 是向量丛 $E$ 在 $x$ 的某个邻域 $U$ 上的截面，使得 $t_1(x), \dots, t_k(x)$ 线性无关。则存在 $x$ 的邻域 $V$，使得 $\forall y \in V$，$t_1(y), \dots, t_k(y)$ 都线性无关。\end{lem}

\begin{proof} 设 $s_1, \dots, s_n$ 是 $x$ 处的一组局部基，并设 $t_i(y) = \sum_j a_{ij}(y)s_j(y)$，其中 $y \mapsto a_{ij}(y)$ 是连续映射 $U \to \GL_n(\C)$。则由假设，$a_{ij}(y)$ 的某个 $k \times k$ 子矩阵必须非奇异，而这对所有充分接近 $x$ 的 $y$ 都应成立。由此即得结论。\end{proof}

% [原文 p.97]
\begin{prop}{2} 设 $f : E \to E'$ 是向量丛映射。则以下条件等价（T.F.A.E.）：
\begin{enumerate}[label=(\arabic*)]
\item $\operatorname{Im}f$ 是 $E'$ 的子丛。
\item $\operatorname{Ker}f$ 是 $E$ 的子丛。
\item $\operatorname{Im}f$ 的纤维的维数局部常值。
\item $\operatorname{Ker}f$ 的纤维的维数局部常值。
\end{enumerate}
\end{prop}

% [Tutor Note: p.97 命题 2 的反例 $f(x, y) = (x, xy)$：在 $x = 0$ 处像的纤维维数从 1 掉到 0，而 $\operatorname{Ker}f$ 在 $x = 0$ 处维数从 0 跳到 1，故两者都不是子丛。原文 "so thy are linearly independent" 中 thy 应为 they。]
\begin{proof} 我们先作一点评论：由线性代数我们知道，对有限维向量空间之间的任意线性映射 $f : V \to V'$，$\operatorname{Ker}f$ 是 $V$ 的向量子空间，$\operatorname{Im}f$ 是 $V'$ 的向量子空间。这不能搬到向量丛上：即若 $f : E \to E'$ 是向量丛映射，则 $\operatorname{Ker}f$ 与 $\operatorname{Im}f$ 不一定是 $E$ 与 $E'$ 的向量子丛。例如，取 $X = I = [0, 1]$ 为单位区间，$E = I \times \C$，$\pi(x, y) = x$。设 $f : E \to E$ 由 $f(x, y) = (x, xy)$ 给出。则 $f$ 的像除 $x = 0$ 处纤维为零外，处处纤维维数为 1，因此 $\operatorname{Im}f$ 不可能是向量丛。$\operatorname{Ker}f$ 也不是。然而这却是唯一可能出错的地方。

回到我们的证明：显然 (3) 与 (4) 等价，且它们由 (1) 或 (2) 推出。为看出 (3) 蕴涵 (1)，取 $x \in X$，选取 $E$ 在 $x$ 处的一组局部基 $s_1, \dots, s_m$ 与 $E'$ 在 $x$ 处的一组局部基 $t_1, \dots, t_n$。设 $k$ 是 $\operatorname{Im}f$ 在 $x$ 处纤维的维数。经过可能的重新编号，我们可以假定 $fs_1(x), \dots, fs_k(x)$ 张成 $\operatorname{Im}f$ 在 $x$ 处的纤维，因而它们线性无关。再经过一次可能的重新编号，我们可以假定 $fs_1(x), \dots, fs_k(x), t_{k+1}(x), \dots, t_n(x)$ 线性无关，于是由 $E'$ 纤维维数的局部常值性与引理 5，可知 $fs_1, \dots, fs_k, t_{k+1}, \dots, t_n$ 构成 $E'$ 在 $x$ 处的一组局部基。由引理 5 的假设，$fs_1, \dots, fs_k$ 构成 $\operatorname{Im}f$ 在 $x$ 处的一组局部基，这意味着 $\operatorname{Im}f$ 是 $E'$ 的子丛。

为看出 (3) $\Rightarrow$ (2)，设 $s_1, \dots, s_m$ 如上。对 $x$ 附近所有 $y$，当 $i > k$ 时可以写 $fs_i(y) = \sum_{j=1}^{k} a_{ij}(y)fs_j(y)$。令 $s'_i(y) = s_i(y) - \sum_{j=1}^{k} a_{ij}(y)s_j(y)$。则 $s'_{k+1}, \dots, s'_m$ 是 $\operatorname{Ker}f$ 的局部截面，且在 $x$ 附近线性无关。由于它们的个数恰好正确，它们构成 $\operatorname{Ker}f$ 的一组局部基，因此 $\operatorname{Ker}f$ 是 $E$ 的子丛。

% [Tutor Note: p.97 段末这条注原文未编号（斜体 "Remark:"），但 p.98 却称引它为 "Remark 1"（见下页批注）；结论 $\operatorname{dim}F_y(\operatorname{Im}f) \ge \operatorname{dim}F_x(\operatorname{Im}f)$ 说明纤维维数是下半连续的：局部只会往上跳，不会往下掉，这正是 (3)$\Rightarrow$(1) 的关键。]
\para{注：} 在没有任何假设的情况下，本证明表明：若 $\operatorname{dim}F_x(\operatorname{Im}f) = n$，则在 $x$ 的某个邻域内对一切 $y$ 有 $\operatorname{dim}F_y(\operatorname{Im}f) \ge n$。\end{proof}

% [原文 p.98]
\begin{proof}[断言 3 的证明] 设 $P$ 是有限生成投射的。则 $P$ 是某个有限生成自由 $C(X)$-模 $F$ 的直和项。因此存在幂等自同态 $g : F \to F$，使 $P \simeq \operatorname{Im}g$。现在 $F = \Gamma(\C^n_X)$ 对某个 $n$ 成立。由上面的命题 1，$g = \Gamma(f)$，其中 $f : \C^n_X \to \C^n_X$。由于 $g^2 = g$，命题 1 蕴涵 $f^2 = f$ 亦然。如果知道 $E = \operatorname{Im}f$ 是 $\C^n_X)$（对某个 $n$，即平凡丛）的子丛，那么由引理 1.1.1 便有 $\C^n_X) \simeq E \oplus E'$，其中 $E' \simeq \operatorname{Ker}(f)$，从而 $P \simeq \operatorname{Im}\Gamma(f) = \Gamma(E)$，因为 $\Gamma(-)$ 是可加函子。由上面的命题 2，只需证明 $\operatorname{dim}F_x(E)$ 局部常值（回忆 $F_x(E)$ 表示 $E$ 在 $x$ 处的纤维）。由于 $f^2 = f$，于是 $E = \operatorname{Ker}f = \operatorname{Im}(1-f)$，且 $F_x(\C^n_X) = F_x(E) \oplus F_x(E')$。设 $\operatorname{dim}F_x(E) = h$、$\operatorname{dim}F_x(E') = k$。对 $f$ 与 $1-f$ 分别应用命题 2 证明末尾的注 1，便得在 $x$ 的某个邻域内对一切 $y$ 有 $\operatorname{dim}F_y(E) \ge h$ 与 $\operatorname{dim}F_y(E') \ge k$。然而
\[ \operatorname{dim}F_y(E) + \operatorname{dim}F_y(E') = h + k \]
是常数，于是平凡丛的维数是局部常值的。\end{proof}

% [Tutor Note: p.98 原文 "Since $f^2 = f$, then $E = \operatorname{Ker}f = \operatorname{Im}(1-f)$" 与同段开头的 $E = \operatorname{Im}f$ 矛盾：对幂等 $f$ 有 $\operatorname{Im}f = \operatorname{Ker}(1-f)$、$\operatorname{Ker}f = \operatorname{Im}(1-f)$，故此处应写成 $E = \operatorname{Ker}(1-f)$（或把 $E'$ 写成 $E' = \operatorname{Ker}f = \operatorname{Im}(1-f)$）。判断依据：同一句随后用 $E' \simeq \operatorname{Ker}(f)$ 且 $F_x(\C^n_X) = F_x(E) \oplus F_x(E')$。译文照原文。]
% [Tutor Note: p.98 两处 "$\C^n_X)$" 多印了右括号（"subbundle of $\C^n_X)$" 与 "Lemma 1.1.1 that $\C^n_X) \simeq E \oplus E'$"），与前面定义 3、p.93 末句的多余括号同属排版笔误；译文照原样保留。]
于是，向量丛与有限生成投射模之间有着完全的类比；这使人们能够把我们在第 1 章（拓扑 K 理论）中遇到的许多概念搬到这个代数框架中（例如定义稳定同构的模等等）。
